\section*{Термоупругая связь}

Рассмотрим линейную термоупругую связь между напряжениями и деформациями:
\begin{equation}
\sigma_{ij}=\lambda\,\delta_{ij}\,\epsilon_{kk}+2\mu\,\epsilon_{ij}-\gamma\,\delta_{ij}\,T.
\end{equation}
Здесь $\lambda,\mu$ --- параметры Ламе, $\epsilon_{ij}$ --- тензор малых деформаций,
$\gamma$ --- коэффициент термоупругой связи, $T$ --- изменение температуры.

\subsection*{1.\;Тензор напряжений $\sigma_{ij}$}
$\sigma_{ij}$ --- компоненты \textbf{тензора напряжений}, то есть внутренних сил в материале,
отнесённых к единице площади. В трёхмерном случае напряжения в точке описываются матрицей $3\times 3$:
\[
\boldsymbol{\sigma}=
\begin{pmatrix}
\sigma_{11} & \sigma_{12} & \sigma_{13}\\
\sigma_{21} & \sigma_{22} & \sigma_{23}\\
\sigma_{31} & \sigma_{32} & \sigma_{33}
\end{pmatrix}.
\]
Диагональные компоненты $\sigma_{11},\sigma_{22},\sigma_{33}$ называются \textbf{нормальными} напряжениями,
а внедиагональные $\sigma_{12},\sigma_{13},\sigma_{23}$ --- \textbf{касательными} (сдвиговыми).
Индексы $i,j$ показывают, какая именно компонента рассматривается.

\subsection*{2.\;Тензор деформаций $\epsilon_{ij}$: определение и смысл}
$\epsilon_{ij}$ --- компоненты \textbf{тензора малых деформаций}. Он определяется через поле перемещений
$\mathbf{u}=(u_1,u_2,u_3)$:
\begin{equation}
\epsilon_{ij}=\frac{1}{2}\left(\frac{\partial u_i}{\partial x_j}+\frac{\partial u_j}{\partial x_i}\right).
\end{equation}

\paragraph{Почему появляется коэффициент $\mathbf{1/2}$.}
Градиент перемещений $\frac{\partial u_i}{\partial x_j}$ содержит одновременно два эффекта:
\begin{itemize}
\item \textbf{деформацию} (изменение расстояний и углов);
\item \textbf{вращение} (поворот тела как жёсткого целого).
\end{itemize}
Чистое вращение не должно считаться деформацией, потому что при вращении расстояния между точками не меняются.
Поэтому берут \textbf{симметричную часть} градиента перемещений. Симметризация даёт сумму двух производных,
а множитель $1/2$ нужен для корректного масштаба деформации.

Из определения сразу следует симметрия:
\[
\epsilon_{ij}=\epsilon_{ji}.
\]

\paragraph{Физический смысл компонент.}
\begin{itemize}
\item $\epsilon_{11},\epsilon_{22},\epsilon_{33}$ описывают относительные удлинения/сжатия вдоль осей.
\item $\epsilon_{12},\epsilon_{13},\epsilon_{23}$ связаны со сдвигом (изменением углов).
\end{itemize}

\subsection*{3.\;След тензора деформаций $\epsilon_{kk}$ и объёмная деформация}
В формуле используется правило Эйнштейна: повторяющийся индекс означает суммирование.
Поэтому
\begin{equation}
\epsilon_{kk}=\epsilon_{11}+\epsilon_{22}+\epsilon_{33}.
\end{equation}
Величина $\epsilon_{kk}$ называется \textbf{следом} тензора деформаций и характеризует \textbf{объёмную деформацию}.
В приближении малых деформаций выполняется:
\[
\frac{\Delta V}{V}\approx \epsilon_{kk},
\]
то есть $\epsilon_{kk}$ примерно равен относительному изменению объёма.

\subsection*{4.\;Дельта Кронекера $\delta_{ij}$}
\[
\delta_{ij}=
\begin{cases}
1, & i=j,\\
0, & i\neq j.
\end{cases}
\]
$\delta_{ij}$ является компонентами единичного тензора (единичной матрицы).
Если некоторый член записан как $\delta_{ij}(\dots)$, то он влияет только на диагональные компоненты:
$\sigma_{11},\sigma_{22},\sigma_{33}$, и не влияет на касательные компоненты.

Физически это означает \textbf{изотропный} (одинаковый во всех направлениях) вклад,
аналогичный гидростатическому давлению.

\subsection*{5.\;Разбор члена $\lambda\,\delta_{ij}\,\epsilon_{kk}$ (объёмный вклад)}
Член
\[
\lambda\,\delta_{ij}\,\epsilon_{kk}
\]
описывает вклад напряжений, связанный с \textbf{изменением объёма}.

\paragraph{Почему именно такой вид.}
\begin{itemize}
\item $\epsilon_{kk}$ измеряет объёмную деформацию: суммарное расширение/сжатие.
\item $\delta_{ij}$ делает вклад одинаковым по всем направлениям, то есть добавляет одинаковое ``давленческое''
напряжение в $\sigma_{11},\sigma_{22},\sigma_{33}$.
\item Параметр $\lambda$ задаёт, насколько сильно материал реагирует на объёмную деформацию.
\end{itemize}

\paragraph{Как выглядит по компонентам.}
Например, для диагональных компонент:
\[
\sigma_{11}\supset \lambda\,\epsilon_{kk},\quad
\sigma_{22}\supset \lambda\,\epsilon_{kk},\quad
\sigma_{33}\supset \lambda\,\epsilon_{kk},
\]
а для касательных ($i\neq j$) этот член равен нулю.

\subsection*{6.\;Разбор члена $2\mu\,\epsilon_{ij}$ (формоизменяющий вклад)}
Член 
\[
2\mu\,\epsilon_{ij}
\]
описывает вклад напряжений, связанный с \textbf{изменением формы} (в том числе со сдвигом).

\paragraph{Физический смысл $\mu$.}
$\mu$ --- модуль сдвига: чем больше $\mu$, тем труднее материалу изменять форму.
Особенно наглядно это видно для касательных компонент. Например:
\[
\sigma_{12}=2\mu\,\epsilon_{12}.
\]

\paragraph{Почему появляется множитель $\mathbf{2}$.}
Он связан с тем, что деформация определена с коэффициентом $1/2$:
\[
\epsilon_{ij}=\frac{1}{2}\left(\frac{\partial u_i}{\partial x_j}+\frac{\partial u_j}{\partial x_i}\right).
\]
При таком определении стандартная линейная упругая связь приводит к множителю $2\mu$ перед $\epsilon_{ij}$.

\subsection*{7.\;Разбор температурного члена $-\gamma\,\delta_{ij}\,T$}
Температурный вклад в формуле:
\[
-\gamma\,\delta_{ij}\,T.
\]

\paragraph{Почему температурный вклад изотропный (через $\delta_{ij}$).}
Равномерный нагрев изотропного материала вызывает стремление расшириться одинаково во все стороны,
то есть температурный эффект не создаёт выделенного направления и действует как ``сферический'' вклад,
поэтому он пропорционален $\delta_{ij}$ и влияет только на нормальные компоненты.

\paragraph{Что означает $T$.}
$T$ --- изменение температуры относительно некоторого исходного состояния (часто $T=\Delta T$).

\paragraph{Что означает $\gamma$.}
$\gamma$ --- коэффициент термоупругой связи: он задаёт, насколько сильно изменение температуры преобразуется
в эквивалентный вклад в нормальные напряжения.

\paragraph{Почему стоит знак минус.}
При нагреве тело \textbf{само} стремится расшириться (свободная термическая деформация).
Если расширению ничего не мешает, то напряжения почти не возникают.
Напряжения появляются тогда, когда свободному расширению мешают закрепления или неоднородность нагрева.
Минус отражает то, что температурное расширение работает как ``встроенная'' деформация, которая вычитается
из механической части.

\subsection*{8.\;Компонентная запись для наглядности}
Для диагональной компоненты:
\[
\sigma_{11}=\lambda(\epsilon_{11}+\epsilon_{22}+\epsilon_{33})+2\mu\,\epsilon_{11}-\gamma T.
\]
Для касательной компоненты:
\[
\sigma_{12}=2\mu\,\epsilon_{12}.
\]
Из этого видно, что температура напрямую не создаёт касательные напряжения: она влияет только на диагональные
компоненты через $\delta_{ij}$.

